\documentclass[10pt,a4paper]{amsart}
\usepackage[margin=19mm]{geometry}
\usepackage{amsmath,amssymb,mathtools}
\usepackage[scr=rsfs,cal=euler]{mathalfa}
\usepackage{xfrac}
\usepackage[unicode,hidelinks,bookmarks=false]{hyperref}
\newcommand{\ie}{i.e.\ }
\newcommand{\cf}{cf.\ }
\newcommand{\resp}{respectively\ }
\newcommand{\Subsection}{\subsection}
\providecommand{\nobreakdash}{\nobreak\hbox{-}}
\setlength{\emergencystretch}{2em}
\multlinegap0pt
\newcommand{\C}{\mathbb{C}}
\newcommand{\Q}{\mathbb{Q}}
\newcommand{\T}{\mathbb{T}}
\newcommand{\Var}{\operatorname{Var}}
\newcommand{\Z}{\mathbb{Z}}
\newcommand{\dd}{\,\mathrm{d}}
\newcommand{\wind}{\operatorname{wind}}
\theoremstyle{plain}
\newtheorem{theorem}{Theorem}[section]
\newtheorem{proposition}[theorem]{Proposition}
\newtheorem{lemma}[theorem]{Lemma}
\newtheorem{corollary}[theorem]{Corollary}
\theoremstyle{remark}
\newtheorem{remark}[theorem]{Remark}
\title[The critical mixed-arithmetic four-point HRT theorem]{The critical mixed-arithmetic four-point HRT theorem}
\author{Vignon Oussa}
\date{Source: arXiv:2609.27970v1 (CC BY 4.0)}
\begin{document}
\maketitle

\noindent\emph{Source and licence.} Vignon Oussa, arXiv:2609.27970v1. The inherited source metadata records CC BY 4.0. The bundled licence notice records the origin of that metadata and the exact contents of this local revision.

\noindent\textit{Editorial scope.} Counted unit: the author's lemma \texttt{lem:winding}, with its complete original proof. This revision restores the author's continued-fraction and Denjoy--Koksma context (source0.tex lines 1028--1059) and the four cited bibliography entries. The statement (1067--1076) and proof (1078--1177) are unchanged. Every mathematical excerpt and bibliography entry below is copied exactly from the bundled source. The published item is preserved separately.

\section*{Author context: irrational returns and Denjoy--Koksma}

For \(\theta\notin\Q\), let
\[
  R_\theta x=x-\theta\pmod 1
\]
act on \(\T\). The rotation \(R_\theta\) is ergodic \cite{Petersen}. If
\(p_n/q_n\) are the continued-fraction convergents of \(\theta\), then
\begin{equation}
  q_n\theta-p_n\longrightarrow0.
  \label{eq:cf-return}
\end{equation}
Thus \(R_\theta^{q_n}\) approaches the identity, and the denominators
\(q_n\) provide the natural times at which an iterated cocycle can be
compared with its initial value.

At these return times, accumulated oscillation is controlled by the
Denjoy--Koksma inequality: for every real-valued periodic function \(g\) of
bounded variation,
\begin{equation}
  \sup_{x\in\T}
  \left|
    \sum_{j=0}^{q_n-1}g(x-j\theta)
    -q_n\int_\T g
  \right|
  \leq \Var(g).
  \label{eq:denjoy-koksma}
\end{equation}
For complex-valued \(g\), we apply the estimate to its real and imaginary
parts. In the winding argument below, it controls the derivative of the
accumulated phase; later, it controls the oscillatory part of a periodic
logarithm. Together with \eqref{eq:cf-return}, it converts recurrence of the
base rotation into rigidity of the cocycle. See
\cite{Khinchin,KuipersNiederreiter,KatokHasselblatt}.

\section{The counted unit and the author's complete proof}

\begin{lemma}[Measurable winding obstruction]
\label{lem:winding}
Let \(\theta\notin\Q\), let \(0\neq h\in L^2(\T)\), and let
\(b\in C^2(\T;\C^\times)\). If
\begin{equation}
  h(x-\theta)=b(x)h(x)
  \label{eq:abstract-cocycle}
\end{equation}
almost everywhere, then \(\wind(b)=0\).
\end{lemma}

\begin{proof}
\textbf{Step 1: Remove the zero set of the transfer function.}
Because \(b\) never vanishes, the zero set of \(h\) is invariant, modulo null
sets, under the irrational rotation. Ergodicity implies that this zero set has
measure zero or one. Since \(h\neq0\), it has measure zero. Therefore
\[
  u(x)=\frac{h(x)}{|h(x)|}
\]
is defined almost everywhere and has modulus one.

\smallskip
\noindent
\textbf{Step 2: Separate the winding from the periodic phase.}
Let \(\psi=b/|b|\). Then
\[
  u(x-\theta)=\psi(x)u(x)
\]
almost everywhere and \(\wind(\psi)=\wind(b)\). If this common winding is
\(d\in\Z\), there is a real-valued periodic \(C^2\) function \(g\) such that
\[
  \psi(x)=e^{2\pi i(dx+g(x))}.
\]

\smallskip
\noindent
\textbf{Step 3: Iterate through continued-fraction returns.}
Iteration through \(q_n\) steps gives
\[
  u(x-q_n\theta)=e^{2\pi i\Phi_n(x)}u(x),
\]
where
\begin{equation}
  \Phi_n(x)
  =dq_nx-d\theta\frac{q_n(q_n-1)}2
   +\sum_{j=0}^{q_n-1}g(x-j\theta).
  \label{eq:Phi-def}
\end{equation}
Since \(|u|=1\),
\[
  |e^{2\pi i\Phi_n(x)}-1|
  =|u(x-q_n\theta)-u(x)|.
\]
By \eqref{eq:cf-return} and translation continuity in \(L^2(\T)\),
\begin{equation}
  e^{2\pi i\Phi_n}\longrightarrow1
  \quad\text{in }L^2(\T).
  \label{eq:phase-to-one}
\end{equation}
In particular,
\begin{equation}
  \int_\T e^{2\pi i\Phi_n(x)}\dd x\longrightarrow1.
  \label{eq:phase-integral-one}
\end{equation}

\smallskip
\noindent
\textbf{Step 4: Exclude nonzero winding by oscillation.}
Assume \(d\neq0\). Differentiating \eqref{eq:Phi-def} gives
\[
  \Phi_n'(x)
  =dq_n+\sum_{j=0}^{q_n-1}g'(x-j\theta).
\]
The function \(g'\) has mean zero and bounded variation. Denjoy--Koksma
therefore yields
\[
  \sup_x\left|
    \sum_{j=0}^{q_n-1}g'(x-j\theta)
  \right|
  \leq\Var(g').
\]
Hence \(\inf_x|\Phi_n'(x)|\geq c q_n\) for all large \(n\), with \(c>0\).
Moreover,
\[
  \|\Phi_n''\|_{L^1}
  \leq q_n\|g''\|_{L^1}.
\]

Although \(\Phi_n\) need not be periodic, it satisfies
\[
  \Phi_n(x+1)=\Phi_n(x)+dq_n,
\]
and \(\Phi_n'\) is periodic. Consequently, the boundary term in integration
by parts vanishes:
\[
  \left[
    \frac{e^{2\pi i\Phi_n(x)}}{2\pi i\Phi_n'(x)}
  \right]_{0}^{1}=0.
\]
It follows that
\[
  \left|
    \int_\T e^{2\pi i\Phi_n(x)}\dd x
  \right|
  \leq
  \frac{\|\Phi_n''\|_{L^1}}
       {2\pi(\inf|\Phi_n'|)^2}
  =O(q_n^{-1}).
\]
This contradicts \eqref{eq:phase-integral-one}. Therefore \(d=0\).
\end{proof}

\begin{thebibliography}{34}

\bibitem{KatokHasselblatt}
A.~Katok and B.~Hasselblatt,
\emph{Introduction to the Modern Theory of Dynamical Systems},
Encyclopedia of Mathematics and its Applications, vol.~54,
Cambridge University Press, Cambridge, 1995.

\bibitem{Khinchin}
A.~Ya. Khinchin,
\emph{Continued Fractions},
Dover, Mineola, New York, 1997.

\bibitem{KuipersNiederreiter}
L.~Kuipers and H.~Niederreiter,
\emph{Uniform Distribution of Sequences},
Wiley, New York, 1974.

\bibitem{Petersen}
K.~Petersen,
\emph{Ergodic Theory},
Cambridge Studies in Advanced Mathematics, vol.~2,
Cambridge University Press, Cambridge, 1983.

\end{thebibliography}

\end{document}
